\DocSection{Discussion and Limitations}
\label{sec:discussion}

The corrected formulation preserves expected-return ordering at the level of
target modes. A finite-temperature draw is nevertheless not an
expected-return-optimal policy in the classical sense; it is a randomized
selection rule whose behavior depends on \(\beta\) and the reference measure.
Those modeling choices must be stated rather than described as data-driven
epistemic uncertainty.

The direct objective in Equation~\eqref{eq:correct-elbo} avoids a nonlinear
transformation of a noisy log density, but it can still have high gradient
variance and a restrictive variational family. A factorized statewise proposal
cannot represent correlations between alternative policy modes. Structured
global latent variables, autoregressive assignments, or retained weighted
policy particles are possible extensions, provided their objectives continue
to target Equation~\eqref{eq:policy-target}.

Finally, simulator access makes the method a form of planning with a generative
model. It should not be described as model-free merely because transition
probabilities are never enumerated.

